% ══════════════════════════════════════════════════════════════════
\section{Introduction}\label{sec:intro}
% ══════════════════════════════════════════════════════════════════

An adaptive stream classifier repairs itself; it does not report.
When an Adaptive Random Forest (ARF)~\cite{gomes_arf_2017} replaces a drifting
subtree, predictive loss returns towards its pre-drift level and nothing else
leaves the model: no change timestamp, no localisation, no record a downstream
system can act on. Deployment separates three objectives that the streaming
literature routinely merges.
\emph{(i) Adaptation} minimises predictive loss; the model is its own consumer.
\emph{(ii) Monitoring and diagnosis} establish whether the environment changed,
when, and where; the consumers are the engineers who must decide whether a shift
is benign seasonality or a broken upstream feature.
\emph{(iii) System alarm} notifies operators, downstream pipelines and auditors.
Objectives (ii) and (iii) survive a model that repairs itself perfectly, because
their consumers are not the model. Regulation states the separation directly:
Article~15(4) of the EU Artificial Intelligence Act~\cite{eu_ai_act_2024}
requires providers of high-risk systems that keep learning after deployment to
eliminate or mitigate the feedback loops that continued learning creates, and
Article~72 requires a documented post-market monitoring system collecting
performance evidence over the system's lifetime. United States banking
supervision requires ongoing monitoring of deployed models under independent
effective challenge~\cite{frb_sr26_2_2026}. A model that silently absorbs a
change satisfies none of this.

The architecture we study---a self-updating model observed by a monitor it does
not control---is the ordinary shape of production monitoring. Managed services
attach drift and model-quality jobs to a deployed endpoint on their own
schedule~\cite{gcp_vertex_monitoring, aws_model_monitor}; open-source libraries
run the same checks in pipelines and CI~\cite{evidently_lib}; NannyML names its
object of observation \emph{the monitored model} and estimates that model's
performance from predictions when labels arrive late or never~\cite{nannyml_pe}.
None of these stacks assumes a frozen model, and none couples the monitor to the
model's update path. River~\cite{montiel_river_2021} and MOA ship both halves and
compose them. We claim no census of deployments; we claim that the composition is
the default design pattern and that regulation now requires the monitoring half
to stay independent of the adaptation half.

Composing the two halves closes a loop. Let $\tau^{*}$ denote the change instant,
$e_t \in \{0,1\}$ the classifier's error indicator at time $t$, and $\Delta e$ the
error excess the change induces before any adaptation. The external monitor
observes $e_t$; the classifier reacts to that same $e_t$ by replacing trees. The
measured signal is endogenous. Fault detection under feedback has studied this
configuration for three decades: a controller that rejects a disturbance also
removes it from the residual, and a residual-based detector loses the evidence it
was calibrated on~\cite{chen_patton_1999, gertler_1998, isermann_2006}. ARF's
internal ADWIN detectors~\cite{bifet_adwin_2007} act as a high-gain,
event-triggered corrector on the error rate; the external cumulative monitor is
the residual test. The change is not undone---the environment has shifted and
stays shifted---the loop has removed its signature.

Adaptation therefore converts a persistent change into a transient one. Write
$\tau_{\mathrm{erase}}$ for the erasure time: the delay from $\tau^{*}$ until
ensemble error returns to within tolerance of its pre-drift level. Any monitor
reading $e_t$ has $\tau_{\mathrm{erase}}$ observations of the change, not the
horizon of the stream. Cumulative-evidence monitors of the CUSUM
family~\cite{page_1954, lorden_1971, moustakides_1986} require a fixed quantity of
evidence, accumulated at a rate set by $\Delta e$, before stopping; when the
window closes first, no alarm fires. Stated this way the effect is a power
deficit under window truncation, and the design problem is a sizing problem: fit
the evidence a monitor requires inside the window the classifier leaves open.
Sequential detection of transient signals~\cite{nikiforov_2012,
tartakovsky_2023_transient} supplies the matching criterion---worst-case
probability of missed detection over a window of finite length at a fixed
false-alarm rate---in place of the F1 scores that dominate drift-detector
benchmarks and that fold calibration and delay into one scalar comparable across
no two detectors.

We call \textbf{blind spot} the region of the design space where the window
closes before the required evidence arrives, and \textbf{starvation} the
mechanism that closes it. Two consequences follow on the same threshold.
Raising the monitor's threshold to control false alarms lengthens the required
accumulation and widens the blind spot; lowering it to fit inside
$\tau_{\mathrm{erase}}$ floods a heteroscedastic stream with alarms. Recall
collapse on stationary streams and precision collapse on non-stationary ones are
the two faces of one calibration, not two phenomena
(Fig.~\ref{fig:ontology}). Our real-data streams exhibit the precision face; the
recall face is established on controlled and semi-synthetic streams, and we
mark the distinction wherever results are reported.

% Ontology of the retained effects -- stream S11-a, thesis v4.
% Supersedes the S5 version on three points: the terminal node is the comparison of the evidence
% budget A against the monitor's own requirement R(D, eps, alpha); the blind spot splits at the
% universal detection floor into an information-theoretic regime and a calibration regime; the
% requirement is shown as a function of the false-alarm level. Retained unchanged: the negative
% feedback loop, the single lambda axis carrying starvation and flooding, the P(X) monitor outside
% the loop. Nodes are placed on an explicit millimetre grid so the picture stays inside the 252 pt
% IEEEtran column; no \ref crosses out of this file, so it compiles against the preamble alone.
% Requires in the preamble:
%   \usepackage{tikz}
%   \usetikzlibrary{arrows.meta,positioning,fit,backgrounds}
\begin{figure}[t]
\centering
\begin{tikzpicture}[
  x=1mm, y=1mm,
  font=\scriptsize,
  box/.style   ={draw, rounded corners=1pt, align=center, inner sep=2pt,
                 text width=34mm, minimum height=6mm},
  loop/.style  ={box, fill=black!5},
  outcome/.style={box, text width=24mm},   % not `out`: /tikz/out is reserved by to[out=..]
  bad/.style   ={box, very thick, text width=24mm},
  ar/.style    ={-{Latex[length=1.4mm]}, thin},
  fb/.style    ={-{Latex[length=1.4mm]}, thin, dashed}
]

\node[box] (chg) at (0,0) {Environment change at $\tau^{*}$};
\node[box] (res) at (-6,-12)
  {Error excess $\Delta e$ on the stream $e_t$\\[1pt]\emph{endogenous residual}};
\node[outcome, text width=20mm] (px) at (28,-12)
  {Monitor on $P(X)$:\\outside the loop};

\node[loop, text width=30mm] (adapt) at (-23,-27)
  {Internal correction:\\one replacement, then learning\\[1pt]
   leaves a budget $A$ over $\tau_{\mathrm{erase}}$};
\node[box, text width=30mm] (mon) at (15,-27)
  {External monitor held at\\false-alarm level $\alpha$\\[1pt]
   requires $R(D,\varepsilon,\alpha)$};

\node[box, text width=44mm] (cmp) at (-4,-43)
  {Evidence available $A$ \emph{vs.} evidence required
   $R(D,\varepsilon,\alpha)$};

\node[bad] (blind) at (-26,-57) {$A < R$\\\textbf{blind spot}};
\node[outcome] (late) at (8,-57) {$A \geq R$\\alarm inside the window};

\node[outcome, text width=32mm] (info) at (-22,-74)
  {$A < A_{\min}(\alpha,\varepsilon)$\\\emph{no} monitor held at this
   level detects};
\node[outcome, text width=32mm] (cal) at (18,-74)
  {$A_{\min} \leq A < R$\\\emph{this} calibration does not;\\another one does};

\node[outcome, text width=40mm] (lam) at (10,-91)
  {One threshold $\lambda$: raising it grows $R$ $\Rightarrow$
   \textbf{starvation}; lowering it to fit $A$ $\Rightarrow$ flooding
   under volatility};

\draw[ar] (chg) -- (res);
\draw[ar] (res.south west) -- (adapt.north);
\draw[ar] (res.south east) -- (mon.north);
\draw[ar] (res.east) -- (px.west);
\draw[fb] (adapt.west) -- ++(-3,0) |- (res.west);
\node[font=\tiny, anchor=south] at (-33,-11.2) {negative feedback};
\draw[ar] (adapt.south) -- (cmp.north west);
\draw[ar] (mon.south)   -- (cmp.north east);
\draw[ar] (cmp.south west) -- (blind.north);
\draw[ar] (cmp.south east) -- (late.north);
\draw[ar] (blind.south west) -- (info.north);
\draw[ar] (blind.south east) -- (cal.north west);
\draw[ar] (cal.south) -- (lam.north);

\end{tikzpicture}
\caption{Relations between the retained effects. The classifier acts by negative
feedback on the very stream the monitor reads, so the evidence reaching the
monitor is a budget $A$ bounded by adaptation rather than by the horizon. The
terminal comparison is $A$ against the evidence $R(D,\varepsilon,\alpha)$ the
monitor's own false-alarm level requires: a blind spot is a deficit between two
quantities in one unit, and it splits at the universal detection floor
$A_{\min}(\alpha,\varepsilon)$. Below the floor no monitor held at that level
detects; above it the deficit belongs to this calibration, and another threshold
or another family meets the same budget. Recall collapse and precision collapse
are two settings of one threshold $\lambda$, not two phenomena. A monitor on the
input distribution is not exposed to the loop, at the cost of blindness to a
change in $P(Y\mid X)$ at fixed $P(X)$.}
\label{fig:ontology}
\end{figure}


\noindent\textbf{Contributions.}
\begin{enumerate}[nosep, label=\textbf{(C\arabic*)}]
  \item \textbf{Monitoring under closed-loop adaptation.} We state the
  concept-drift monitoring problem as fault detection on an endogenous residual,
  give the mapping to the fault-detection-under-feedback setting, and derive the
  prediction that severity is governed by the ratio of erasure time to detector
  integration time---not by River, by ARF, or by ADWIN
  (Section~\ref{sec:blindspot}).
  \item \textbf{A transient-detection analysis of starvation.} We characterise
  the observability window $\tau_{\mathrm{erase}}$ induced by adaptation and give
  a finite-horizon sizing condition relating evidence required to evidence
  available, together with a detection floor that no threshold choice escapes
  (Section~\ref{sec:starvation}).
  \item \textbf{Synchronised measurement.} We instrument tree replacement,
  ensemble error and the external monitor's statistic on a shared clock, and
  report how strongly replacement counts predict error recovery, so that erasure
  is measured rather than assumed (Section~\ref{sec:experiments}).
  \item \textbf{An architectural consequence.} We order monitor families by their
  exposure to the loop---cumulative on the error stream, windowed on the error
  stream, distributional on the input space---and state the cost each pays
  (Section~\ref{sec:related} and Section~\ref{sec:discussion}).
\end{enumerate}